% ============================================================
% STRUCTURE AND REFERENCING GUIDE
% Book I: Ontology of Viability
% ============================================================

\cleardoublepage

\chapter*{Structure and Referencing Guide}

\phantomsection
\addcontentsline{toc}{chapter}{Structure and Referencing Guide}

% ------------------------------------------------------------
% 1. PURPOSE OF THIS GUIDE
% ------------------------------------------------------------

\section*{1. Purpose of This Guide}

This guide establishes the conventions used to identify,
locate, reference, and relate the structural and logical
components of \textit{Ontology of Viability}.

The book contains several kinds of identifiable elements.
These include structural units such as parts, chapters,
sections, subsections, and paragraphs; logical entities
such as axioms, definitions, propositions, theories,
hypotheses, arguments, and philosophical frameworks;
internal references connecting different locations within
the book; and external sources used in the development
or evaluation of its arguments.

These elements do not all belong to the same category.
A structural location identifies where something occurs
within the organization of the book. A logical entity
identifier identifies a particular conceptual or
argumentative entity. A reference establishes a relation
between identifiable elements, while a bibliographic
reference identifies an external source.

The purpose of the referencing system is to preserve
these distinctions and make the internal organization
of the book traceable.

The conventions presented in this guide are intended
to support:

\begin{itemize}
    \item stable identification of structural locations;
    \item stable identification of logical entities;
    \item precise internal cross-referencing;
    \item reliable location of entities within the text;
    \item explicit representation of logical dependencies;
    \item distinction between internal entities and
    external sources;
    \item consistent organization of bibliographic
    information.
\end{itemize}

The guide is therefore not an additional theoretical
part of the ontology. It provides the structural and
referential conventions through which the arguments
and conceptual entities of the book can be located
and examined.

% ------------------------------------------------------------
% 2. THE DISTINCTION BETWEEN STRUCTURAL AND LOGICAL
% IDENTIFICATION
% ------------------------------------------------------------

\section*{2. Structural and Logical Identification}

The referencing system distinguishes two principal
forms of identification: structural identification
and logical entity identification.

Structural identification refers to the position of
a textual element within the architecture of the book.
It makes it possible to locate a particular part,
chapter, section, subsection, or paragraph independently
of the conceptual content expressed there.

Logical entity identification, by contrast, refers
to a specific conceptual or argumentative entity.
An entity may be defined, discussed, developed,
criticized, or referenced at more than one structural
location. Its logical identity should therefore not
be confused with the location of any particular
discussion of it.

A structural identifier answers the question:

\begin{quote}
Where is this element located within the structure
of the book?
\end{quote}

A logical entity identifier answers a different question:

\begin{quote}
Which identifiable conceptual or argumentative
entity is being discussed?
\end{quote}

These two forms of identification are complementary.
A logical entity may have one stable identifier while
appearing in several different structural locations.
Conversely, a structural location may contain several
logical entities.

The system therefore does not use structural identifiers
as substitutes for logical entity identifiers, or
logical entity identifiers as substitutes for
structural locations.

% ------------------------------------------------------------
% 3. THE ORGANIZATION OF THIS GUIDE
% ------------------------------------------------------------

\section*{3. The Organization of This Guide}

This guide is divided into seven complementary
documents. Each document addresses a distinct
component of the structural and referencing system.

\subsection*{3.1. Structural Identifiers}

The document \textit{Structural Identifiers}
establishes the convention for identifying the
hierarchical locations of textual elements within
the book.

It explains the structure and use of identifiers
that correspond to the organization of the work,
including the relevant levels of textual division.

The structural identifier system provides a
consistent means of locating passages independently
of page numbering.

\subsection*{3.2. Logical Entity Identifiers}

The document \textit{Logical Entity Identifiers}
establishes the convention for assigning stable
identifiers to logical entities.

The system includes distinct identifier families,
including:

\begin{itemize}
    \item \texttt{AX} --- AX-series identifiers;
    \item \texttt{DF} --- DF-series identifiers;
    \item \texttt{PR} --- PR-series identifiers;
    \item \texttt{TH} --- TH-series identifiers;
    \item \texttt{HY} --- HY-series identifiers;
    \item \texttt{AR} --- AR-series identifiers;
    \item \texttt{PF} --- PF-series identifiers.
\end{itemize}

Examples of the identifier format include
\texttt{AX-001}, \texttt{DF-001}, and
\texttt{PR-001}.

The purpose of these identifiers is to preserve
the identity of a logical entity across different
locations and discussions within the system.

\subsection*{3.3. Internal Cross-References}

The document \textit{Internal Cross-References}
establishes how one structural location or logical
entity may refer to another element within the
book.

It provides the conventions for making internal
references precise, consistent, and traceable.

Internal cross-references should make clear whether
the reference concerns a textual location, a logical
entity, or a relation between them.

\subsection*{3.4. Entity Location References}

The document \textit{Entity Location References}
explains how the structural locations associated
with logical entities are recorded and referenced.

Because a logical entity may be discussed in
multiple places, its identifier alone does not
necessarily indicate every location at which it
appears.

The location-reference convention makes it
possible to distinguish the identity of an entity
from the places where it is introduced, defined,
developed, used, or examined.

\subsection*{3.5. Logical Dependencies}

The document \textit{Logical Dependencies}
establishes the conventions for representing
dependencies among logical entities and
argumentative components.

A proposition, argument, theory, or other
logical entity may depend on definitions,
assumptions, principles, or previously established
claims.

The purpose of this document is to make such
dependencies explicit and traceable without
confusing a logical relation with a mere textual
reference.

\subsection*{3.6. External Source References}

The document \textit{External Source References}
establishes the conventions for referring to
materials that originate outside the internal
logical structure of the Viability Philosophy
System.

External sources must remain distinguishable
from the system's own logical entities.

An external source does not acquire an internal
logical entity identifier merely because it is
quoted, discussed, cited, or used as a point of
comparison.

The identification and referencing of external
materials must therefore follow the conventions
specified in the corresponding document.

\subsection*{3.7. Bibliographic Organization}

The document \textit{Bibliographic Organization}
establishes the conventions for organizing
bibliographic information and maintaining
consistency between references in the text
and bibliographic entries.

Its purpose is to provide a systematic
bibliographic structure while preserving
the distinction between the internal
identification system of the book and
the bibliographic identification of
external publications.

% ------------------------------------------------------------
% 4. GENERAL PRINCIPLES
% ------------------------------------------------------------

\section*{4. General Principles}

The structural and referencing system is governed
by several general principles.

\subsection*{4.1. Stability}

An identifier should remain stable when the
corresponding entity or structural element
continues to be the same identifiable element.

Changes in wording, exposition, or textual
location do not automatically imply that a
logical entity has acquired a new identity.

Likewise, the assignment of an identifier
should not be changed casually once it has
become part of the internal referencing system.

\subsection*{4.2. Distinction of Categories}

Different kinds of elements must be identified
according to their respective categories.

A structural location, a logical entity,
an internal reference, and an external
bibliographic source are not interchangeable
objects of identification.

The system must preserve these distinctions
even when several categories are involved
in the same discussion.

\subsection*{4.3. Traceability}

A reference should make it possible to
identify the element to which it points.

Where appropriate, a logical entity should
be traceable to its relevant structural
locations, and a structural location should
be identifiable within the hierarchy of
the book.

Traceability is necessary for examining
the development of an argument and
following the relations among its
conceptual components.

\subsection*{4.4. Consistency}

The same conventions should be applied
consistently throughout the book.

An identifier should not be used to
designate different entities, and
different identifiers should not be
assigned to the same entity without
a clearly established reason.

The rules for structural identification,
logical entity identification, internal
references, and bibliographic organization
must be applied according to their
respective documents.

\subsection*{4.5. Separation of Identity and Location}

The identity of a logical entity must
remain distinct from its location in
the text.

An entity may be discussed in more
than one place, and its structural
locations may change as the manuscript
is revised. Such changes should not
automatically destroy the continuity
of its logical identification.

This principle is particularly important
for maintaining stable references during
the development and revision of the book.

\subsection*{4.6. Separation of Internal and External
Identification}

The internal identification system
belongs to the architecture of the
Viability Philosophy System.

External publications and other
external materials retain their own
bibliographic identities and are
referenced according to the relevant
external-source conventions.

An external source should not be
assigned an internal logical entity
identifier merely because it contributes
to an argument or is discussed in the text.

% ------------------------------------------------------------
% 5. USE OF THE GUIDE
% ------------------------------------------------------------

\section*{5. Use of the Guide}

The seven documents that follow should
be consulted according to the type of
identification or reference required.

For the identification of a textual
location, consult \textit{Structural
Identifiers}.

For the identification of a conceptual
or argumentative entity, consult
\textit{Logical Entity Identifiers}.

For references connecting locations
or entities within the book, consult
\textit{Internal Cross-References}.

For recording the places where a
logical entity occurs, consult
\textit{Entity Location References}.

For representing relations of
dependence among logical entities,
consult \textit{Logical Dependencies}.

For references to materials originating
outside the system, consult
\textit{External Source References}.

For the organization of bibliographic
entries and their relation to textual
citations, consult
\textit{Bibliographic Organization}.

These documents are complementary.
No single convention should be used
as a substitute for another when
different kinds of identification
are required.

The guide should be read together
with the relevant conventions in
the subsequent documents. In cases
where a specific rule is established
in one of those documents, that
document governs the corresponding
technical procedure.

% ------------------------------------------------------------
% 6. THE FUNCTION OF THE REFERENCING SYSTEM
% ------------------------------------------------------------

\section*{6. The Function of the Referencing System}

The referencing system is intended
to support the systematic examination
of the book's conceptual and argumentative
structure.

Its function is not merely to facilitate
navigation through the text. It also
makes it possible to distinguish the
identity of a claim from the location
at which it is presented, to follow
the development of a logical entity,
and to examine the dependencies and
relations among different components
of the argument.

The system also supports the distinction
between the internal conceptual
architecture of the book and the
external materials with which it
engages.

In this way, structural identification,
logical entity identification, internal
referencing, dependency representation,
and bibliographic organization contribute
to a unified framework of scholarly
traceability.

The purpose of this framework is to
make the organization of the work
explicit, preserve the continuity
of its identifiable components, and
enable its arguments to be examined
with precision.

\cleardoublepage

% ============================================================
% END OF STRUCTURE AND REFERENCING GUIDE
% ============================================================